\section{A Method for the Next Decision}

The book began with a mismatch: proposals appear at inference speed while the world changes through many clocks. The answer is a way to inspect the route between them.

\subsection{Ask What the Capability Establishes}

A convincing interface can establish useful substitutability for a specified task. It does not settle every question about the systems behind it.

Turing's behavioral approach helped make machine intelligence an operational subject. A test's result still belongs to its test conditions. Producing an answer, maintaining a safe service, possessing a right, and bearing an obligation use different criteria.

A person's worth consequently cannot be reduced to the current scarcity of a particular output. Capabilities matter because they expand what a person and society can do. Their changing market price does not exhaust the value of a life.

Experience, relationship, participation, care, and judgment remain dimensions on which a society can choose to use its growing capacity. They need institutions and resources, rather than a promise that displaced output will automatically reappear in a higher form.

\subsection{Recalculate the Scarce Interface}

When a new tool lowers a cost, inspect what becomes scarce next.

Candidate production may give way to review. Review may improve through automatic checking and reveal a specification bottleneck. Deployment may reveal maintenance. Widely trusted service may reveal an access or distribution problem.

The movement has no fixed occupational destination. A slow activity can be improved. A trusted institution can become ceremonial. A supposedly protected skill can become easy to assess or automate.

The enduring questions concern the function. What result does this position establish? Which evidence supports it? Could another arrangement supply the same function better? Who controls access to that arrangement?

The relative importance of a function does not guarantee a good livelihood for its worker. Pay, training access, ownership, and bargaining power need their own account. Otherwise a prediction about growing care needs becomes merely a prediction about growing unpaid work.

\subsection{Three Questions Before Removing Friction}

\textbf{What does the change actually remove?} Specify the task, baseline, quality, time, and system boundary. A lower drafting time is a real gain, while a lower time to a dependable decision is a different claim.

\textbf{What happens elsewhere?} Examine review, rework, energy, infrastructure, support, exposure, and recovery. Some costs move; some fall; some appear for the first time. Measure the net result instead of assuming conservation.

\textbf{Who receives the benefit and the remaining bill?} Identify users, workers, operators, affected third parties, and future obligations. A cheaper interface can have an unequal distribution of gains even when total cost falls.

For a software team, these questions connect generation to review, regression risk, and maintenance. For a clinical service, they connect a model's measured performance to intended use, professional duties, patient rights, and follow-up. For a fleet, they connect vehicle efficiency to empty miles, coverage, infrastructure, and accident response.

The answers should be concrete enough to change the proposal. If the project discovers an overloaded review stage, it can improve selection or capacity. If a failure cannot be contained, it can narrow scope or redesign recovery. If the bill has no responsible bearer, it needs a responsibility arrangement before broad deployment.

\subsection{Keep a Few Capabilities in Practice}

At the individual level, preserve the abilities on which your own judgment depends. Let the tool handle what it does well; retain enough direct work to understand, inspect, and revise the result you rely on.

The amount differs by task. An occasional practice session can preserve one skill. Another requires continued supervised exposure. Good assistance can make both more effective.

At the organizational level, maintain the ability to reconstruct a decision. Keep proportionate records, competent reviewers, independent challenge, and usable authority to stop or correct the process.

These functions need a budget and a role. An instruction to “be careful” cannot substitute for time, access, training, and power.

Revisit retained friction too. A manual step once needed for safety may become replaceable by a better check. A learning exercise may cease to teach. Preserving the purpose requires changing the method when evidence supports change.

\subsection{Protect Choice Where Costs Are Imposed}

People affected by a decision need information and effective representation. Some can consent directly. Others require professional duties, public procedures, guardianship, or protections for interests that cannot speak for themselves.

Authority should fit competence and control, while remedies should remain usable by those carrying the loss. These are design criteria for a legitimate arrangement, not a claim that every affected person must personally operate every system.

The method can be repeated as capability grows. Identify the receiving medium. Separate the constraints. Trace the conversion into action. Examine the feedback and the bill.

A better tool changes the answers. It does not make the questions obsolete.
